% moved from main.tex to the supplementary material (2026-11-09)
\begin{table*}[t]
\caption{Departures from the original implementations.}
\label{tab:deviations}
\small
\begin{adjustbox}{max width=\linewidth}\begin{tabular}{p{1.8cm}p{14.3cm}}
\toprule
Technique & Departure and reason \\
\midrule
BLAZE & Hard examples are weighted inside the batch by the loss; the paper's GPT-mined hard-example set is BeetleBox-specific. Comments are not stripped before chunking. The learning rate is selected by validation MAP from $\{2{\times}10^{-7}, 2{\times}10^{-6}, 2{\times}10^{-5}\}$ (the paper's value and two larger ones), because the paper's value was set for $\sim$11{,}000 training bugs and our projects have 100--600. Selected in fold~1: $2{\times}10^{-7}$ for liquibase and localstack, $2{\times}10^{-6}$ for dolphinscheduler, $2{\times}10^{-5}$ for elasticsearch, ansible and airflow; all selected values in the learning-rate table of the replication package. \\
FLIM & Positive pairs are the functions the fix changed (diff overlap), not every function of the fixed file. Its TF-IDF features split identifiers at camel-case boundaries without keeping the whole identifier. The CodeSearchNet pairs of the paper's MIX set are not added. IR feature f2 (API-enriched similarity) needs JDK documentation and has no Python counterpart, so it is dropped. Other positives of the same report are masked out of the in-batch loss. The learning rate is selected by validation MAP between the two published values ($2{\times}10^{-5}$ in the released code, $5{\times}10^{-5}$ in the paper). Selected in fold~1: $2{\times}10^{-5}$ for elasticsearch and airflow, $5{\times}10^{-5}$ for the other four projects; the learning-rate table of the replication package. \\
LocAgent & No change to prompts, tools or search logic. One bug in its driver is fixed: a tool output is passed to \texttt{eval()}, which fell back to the raw text only on a syntax error, so outputs beginning with a path such as \texttt{lib/ansible/\ldots} raised a name error, killed the worker and left the run waiting forever; the fix extends the existing fallback to every error, so every output that did not crash is handled exactly as before. An adapter exports our test bugs as LocAgent instances (pre-fix commit, report text; the patch field is left empty) and builds LocAgent's graph and BM25 indexes from our local clones with LocAgent's own builders, materialising the full commit tree (not \texttt{git archive}, which drops paths marked \texttt{export-ignore}), so LocAgent sees the same pre-fix snapshot as the other techniques. We use the paper's fine-tuned 7B model rather than its 32B model: the 32B bf16 weights (65.5\,GB) exceed our 48\,GB GPU, and loading them with on-the-fly FP8 quantization ran out of memory. The paper reports both; its ablations use the 7B model (file-level Acc@5 on SWE-bench Lite 88.3 for 7B, 92.7 for 32B). \texttt{litellm} 1.52.1, pinned by LocAgent, is no longer on PyPI; 1.53.1 is used. \\
SweRank & No change to the corpus builder, retriever, prompts, reranking loop or permutation parsing; reranking uses the authors' sequential path, one bug at a time. An adapter feeds the same bugs, pre-fix snapshots and report text as for LocAgent, and maps the reranked functions to files by first appearance; files beyond the top 100 functions keep the retriever's order. SweRank's structure builder files a snapshot's top-level files under the repository name, which collides with a package directory of the same name (\texttt{localstack/}, \texttt{airflow/}); file paths are resolved against the snapshot. A function text already embedded for an earlier bug of the same project is not embedded again. vLLM's memory share is lowered from 0.9 to 0.55 for our shared GPU. We use the Small models (7B reranker), a configuration documented in the released code, rather than the Large ones: the 32B reranker exceeds our GPU, as for LocAgent. A faster variant that sent the reranking prompts of many bugs to the GPU together did not reproduce the sequential order (0 of 6 bugs) and was not used. \texttt{pyarrow}, which the authors do not pin, is set to 20.0.0. \\
\bottomrule
\end{tabular}\end{adjustbox}
\end{table*}
